% Généré par pertesFer/calculPertes/resultatsNote.m le 17/11/2026 16:30. Ne pas éditer.
\newcommand{\vitRef}{1500}
\newcommand{\tauCons}{2}
\newcommand{\jeuCarte}{m330\_35a}
\newcommand{\nPositions}{90}
\begin{table}[H]
\centering
\caption{Trajectoires simulées avec FEMM~5 (90 positions sur un pas polaire). $I_s$ : amplitude du courant.}
\label{tab-simulations}
\resizebox{\textwidth}{!}{%
\begin{tabular}{llS[table-format=2.2]S[table-format=2.2]S[table-format=1.3]S[table-format=1.3]S[table-format=5.0]S[table-format=5.0]}
\toprule
Machine & Loi & {$I_s$ moyen (A)} & {$I_s$ max (A)} & {Couple moyen (N\,m)} & {Ondulation c-c (N\,m)} & {Éléments stator} & {Éléments rotor} \\
\midrule
sans aimant & P0 & 9.56 & 9.56 & 2.000 & 1.035 & 5107 & 8537 \\
sans aimant & P1 & 9.63 & 10.84 & 2.001 & 0.013 & 5107 & 8537 \\
avec aimants (B123\_f25) & P0 & 6.70 & 6.70 & 2.000 & 0.965 & 5107 & 8820 \\
avec aimants (B123\_f25) & P1 & 6.73 & 7.84 & 2.001 & 0.009 & 5107 & 8820 \\
\bottomrule
\end{tabular}}
\end{table}

\begin{table}[H]
\centering
\caption{Pertes fer à \SI{1500}{tr/min} et \SI{2}{\newton\metre}, machine entière (W), pour chaque jeu de coefficients.}
\label{tab-pertes}
\resizebox{\textwidth}{!}{%
\begin{tabular}{llS[table-format=2.2]S[table-format=2.2]S[table-format=1.2]S[table-format=2.2]S[table-format=2.2]S[table-format=2.2]}
\toprule
Machine & Loi & {Hystérésis} & {Classiques} & {Excédentaires} & {Stator} & {Rotor} & {Total} \\
\midrule
\multicolumn{8}{l}{\emph{Jeu « m330\_35a »}} \\
sans aimant & P0 & 5.85 & 2.03 & 2.90 & 7.70 & 3.08 & 10.78 \\
sans aimant & P1 & 6.43 & 2.65 & 3.52 & 9.14 & 3.46 & 12.60 \\
avec aimants (B123\_f25) & P0 & 4.14 & 1.29 & 2.04 & 5.75 & 1.73 & 7.48 \\
avec aimants (B123\_f25) & P1 & 4.83 & 2.27 & 3.15 & 7.88 & 2.36 & 10.24 \\
\midrule
\multicolumn{8}{l}{\emph{Jeu « initial »}} \\
sans aimant & P0 & 7.18 & 1.70 & 0.00 & 7.49 & 1.40 & 8.89 \\
sans aimant & P1 & 8.02 & 2.23 & 0.00 & 8.66 & 1.59 & 10.25 \\
avec aimants (B123\_f25) & P0 & 4.79 & 1.09 & 0.00 & 5.16 & 0.73 & 5.88 \\
avec aimants (B123\_f25) & P1 & 5.65 & 1.91 & 0.00 & 6.55 & 1.01 & 7.56 \\
\midrule
\multicolumn{8}{l}{\emph{Jeu « hugo »}} \\
sans aimant & P0 & 5.26 & 7.95 & 1.56 & 9.42 & 5.34 & 14.77 \\
sans aimant & P1 & 5.76 & 10.40 & 1.90 & 11.93 & 6.13 & 18.06 \\
avec aimants (B123\_f25) & P0 & 3.74 & 5.07 & 1.10 & 7.02 & 2.90 & 9.92 \\
avec aimants (B123\_f25) & P1 & 4.36 & 8.91 & 1.70 & 10.83 & 4.14 & 14.97 \\
\midrule
\multicolumn{8}{l}{\emph{Jeu « boglietti3335 »}} \\
sans aimant & P0 & 4.65 & 3.14 & 0.00 & 5.77 & 2.02 & 7.79 \\
sans aimant & P1 & 5.15 & 4.11 & 0.00 & 6.93 & 2.32 & 9.26 \\
avec aimants (B123\_f25) & P0 & 3.20 & 2.00 & 0.00 & 4.14 & 1.07 & 5.21 \\
avec aimants (B123\_f25) & P1 & 3.75 & 3.52 & 0.00 & 5.74 & 1.52 & 7.26 \\
\bottomrule
\end{tabular}}
\end{table}

\begin{table}[H]
\centering
\caption{Pertes par hystérésis à \SI{1500}{tr/min} (W), jeu « m330\_35a » : comptage des cycles (formule retenue) et somme sur les harmoniques (ancienne écriture).}
\label{tab-hysteresis}
\resizebox{\textwidth}{!}{%
\begin{tabular}{llS[table-format=1.2]S[table-format=1.2]S[table-format=1.2]S[table-format=1.2]}
\toprule
Machine & Loi & {Stator, cycles} & {Stator, harmoniques} & {Rotor, cycles} & {Rotor, harmoniques} \\
\midrule
sans aimant & P0 & 5.03 & 5.11 & 0.82 & 0.84 \\
sans aimant & P1 & 5.53 & 5.37 & 0.89 & 0.93 \\
avec aimants (B123\_f25) & P0 & 3.69 & 3.69 & 0.45 & 0.46 \\
avec aimants (B123\_f25) & P1 & 4.27 & 4.08 & 0.56 & 0.58 \\
\bottomrule
\end{tabular}}
\end{table}
